% 可直接并入总论文。依据 origin/main 的 b3c23ed 版本及其中四问最终检查说明撰写。
\section{模型评价、改进与推广}

\subsection{模型评价}
本文将多尺度因素识别、时空风险评价、滚动触发预报与二维传播模拟相衔接，既能解释灾前环境异常，也能给出危险窗口、口岸到达时间和条件影响范围。各模块保留温度、辐射、水分、地形及物源等物理含义，便于追溯风险变化的来源。因素排序采用年份整块自助法与相关去冗余检验，空间评价采用分块交叉验证，传播模型另设下游留出区检验；集合情景则将天气和参数的不确定性传递至最终结果，使模型评价同时兼顾点估计与区间表现。

模型目前主要受同机理事件数量和源区观测条件限制。历史背景年份较少，风险阈值随年份选择存在波动；远期起报的触发窗口较宽，传播结果也受地形分辨率、入沟体积及摩阻参数影响。因此，现阶段结果适合吉隆本地事件分析与条件情景评价，跨事件的命中率和虚警率仍需长期资料检验。

\subsection{改进方向}
后续应优先补充同机理事件及长期业务预报档案，在独立事件上检验触发阈值、预警提前时间和区间覆盖情况。源区可增加气温、融水与冰岩形变监测，减小区域观测向高海拔源区迁移的误差；传播计算可结合更精细的灾前地形、物源调查及沿程流深、流速观测，联合约束有效入沟体积、摩阻和侵蚀卷入参数。滚动预报继续区分严格预测与本地回放校准，使新增资料用于检验实际预警能力。

\subsection{推广应用}
这一框架可用于同类高山冰岩崩塌灾害链的风险筛查、重点沟谷监测布设及应急疏散情景比选。迁移到其他区域时，可保留因素识别、条件触发与传播计算的结构，依据当地气候、物源和沟道条件重新建立背景分布与参数范围。对于降雨主导的泥石流，应以降雨—水分过程为主要触发分支，重新标定入口过程和流变参数；相对风险等级与条件影响区可作为分级巡查和避险路线选择的辅助依据。
